% ============================================
% Economist Cover Letter – U.S. Department of the Treasury
% ============================================

\documentclass[11pt,letterpaper]{letter}

\usepackage{geometry}
\usepackage{microtype}
\usepackage[hidelinks]{hyperref}

\geometry{left=25mm,right=25mm,top=25mm,bottom=25mm}

\signature{Decory Edwards}
\address{
Department of Economics \\
Johns Hopkins University \\
Baltimore, MD 21218 \\
\texttt{dedwar65@jh.edu}
}

\begin{document}

\begin{letter}{Hiring Committee \\
Office of the Assistant Secretary for Economic Policy \\
U.S. Department of the Treasury \\
Washington, DC}

\opening{Dear Hiring Committee,}

I am writing to apply for the Economist position in Macroeconomic Analysis within the Office of the Assistant Secretary for Economic Policy at the U.S. Department of the Treasury. I am a Ph.D.\ candidate in Economics at Johns Hopkins University, advised by Professors Christopher D.~Carroll, Nicholas W.~Papageorge, and Stelios Fourakis, and I expect to receive my degree in May 2026. My research fields are macroeconomics, wealth and income dynamics, and computational economics.

My research agenda focuses on understanding the determinants of wealth inequality and the mechanisms through which heterogeneity in returns to wealth shapes aggregate outcomes. In my job-market paper, I estimate the degree of heterogeneity in individual portfolio returns using micro-data from the Survey of Consumer Finances and simulate a structural model in which households face idiosyncratic return risk. I show that allowing for persistent heterogeneity in returns substantially improves the model’s ability to match the empirical distribution of wealth, offering a tractable way to connect micro-level return dispersion to macroeconomic inequality.

In a second project, I use the Health and Retirement Study to examine the relationship between interpersonal trust and portfolio performance. Building on the literature that finds a hump-shaped relationship between trust and income, I show that a similar relationship holds for individual returns to wealth measured in the HRS data, highlighting a behavioral channel through which social attitudes shape saving and investment outcomes. This work also provides new estimates of the persistent component of returns to net worth, which motivate the calibration of heterogeneity in my structural model.

I am particularly interested in the Office of Economic Policy because of its role in producing policy-relevant macroeconomic analysis that directly informs fiscal decisions, economic forecasting, and public communication. Treasury’s work sits at the intersection of data, theory, and real-time policy needs. This environment closely matches my training in combining micro-level evidence with structural macroeconomic frameworks to evaluate aggregate dynamics and distributional consequences of policy.

My research has direct relevance for several core areas of Treasury’s macroeconomic work, including consumption dynamics, wealth concentration, and the transmission of fiscal and monetary policy across households. I am comfortable working with large administrative and survey datasets, producing transparent empirical analysis under time constraints, and translating technical results into clear economic narratives that can inform decision-makers. I am especially motivated by the opportunity to contribute analysis that supports public service and improves understanding of how macroeconomic policy affects households across the income and wealth distribution.

I would welcome the opportunity to contribute to the Macroeconomic Analysis team and to apply my quantitative and empirical training in a policy-focused setting. Thank you very much for your time and consideration. I would be glad to provide any additional information and look forward to the possibility of discussing my application.

\closing{Sincerely,}

\end{letter}
\end{document}
